\documentclass[twocolumn]{aastex631}
\usepackage{amsmath}
\usepackage{hyperref}

\shorttitle{Symbolic Regression and the Hubble Constant}
\shortauthors{Hernandez}

\begin{document}

\title{The Hubble Constant from Symbolic Regression: DESI DR2 Confirmation,
M(z) Non-Evolution, and External Constraints}

\author{Ivan Hernandez}
\affiliation{Independent Researcher}

\begin{abstract}
We apply symbolic regression via PySR to cosmic chronometer (CC), BAO (SDSS +
DESI DR2), and Pantheon+ supernova data to discover the cosmic expansion
history $H(z)$ with minimal theoretical priors.  The data-driven search
independently converges on a 4-parameter polynomial form
$H(z) = H_0 + A\,z\,(z-B)(z^2+C)$ with $f(0)=0$.  Joint profile likelihood
yields $H_0 = 68.0\ [67.3, 69.1]$ km/s/Mpc (68\% CL), consistent with Planck 2018 at
$1.2\sigma$.  This result is robust
across: three independent supernova samples (Pantheon+, DES-SN5YR, Union3);
three BAO configurations (SDSS, DESI DR1, DESI DR2); seven systematic
objection tests, and an adversarial validation exercise (14 challenges across
2 rounds, none fatal).  DESI DR2 confirms the DR1 result with
factor-two improved precision.
We test for evolution of the supernova absolute magnitude with redshift,
finding $\alpha = 0.020 \pm 0.010$ (68\% CL), consistent with zero.
External constraints --- GW170817 (65.5$\pm$4.4), DES Y3+GW (67.9$\pm$4.4),
and TDCOSMO 2025 lensing (71.6$\pm$3.6) --- combine to $H_0=68.8\pm2.3$,
consistent with our result.  The lone test that changes $H_0$ is fixing the
absolute magnitude $M$ to the SH0ES Cepheid calibration, which forces
$H_0=73.7$ with $\Delta\chi^2_{\rm SN}=+50$ ($7.1\sigma$ rejection).
A direct flat $\Lambda$CDM fit to the same data gives $H_0=68.6$,
$\Omega_m=0.306$, joint
$\chi^2=1430.4$ --- nearly identical to the SR solution ($\chi^2=1430.0$),
confirming that the data-driven expansion history is consistent with the
standard model without assuming it.  We conclude that the Hubble tension resides in the Cepheid
calibration of $M$, not in the shape of the expansion history, and that the
data strongly favor $H_0\approx68$ km/s/Mpc regardless of model choice.
\end{abstract}

\section{Introduction}

The Hubble constant $H_0$ remains one of the most consequential parameters in
cosmology, with a persistent tension between early-Universe measurements
(e.g., Planck CMB~\cite{planck2018}: $H_0 = 67.4 \pm 0.5$) and the local
distance-ladder measurement by SH0ES~\cite{sh0es2024}: $H_0 = 73.0 \pm 1.0$).
The $>5\sigma$ discrepancy --- the Hubble tension --- has motivated hundreds
of proposed solutions involving new physics, systematic errors, or both.

Most analyses assume a specific cosmological model (typically flat
$\Lambda$CDM) to constrain $H_0$ from expansion-history data.  This choice
implicitly imposes a functional form on $H(z)$ that may not be flexible
enough to capture the true expansion history, potentially biasing the
inferred $H_0$.

Here we take a different approach: we use symbolic regression (PySR~\cite{pysr})
to discover the functional form of $H(z)$ directly from the data, with
minimal priors.  We then use the discovered form --- independently verified
across 6 pre-DESI SR seeds (all converging to the same 4-parameter cubic
family) and confirmed by 2 DESI-inclusive runs --- to jointly fit cosmic chronometer,
BAO, and supernova data, marginalizing over the supernova absolute magnitude
$M$ to avoid any Cepheid calibration bias.

\section{Data}

\subsection{Cosmic Chronometers (33 points)}
We use 33 $H(z)$ measurements from cosmic chronometry \cite{cc_data},
spanning $z\in[0.07,1.965]$.  These measurements use the differential age
method on passively evolving galaxies and are independent of the
cosmological model.

\subsection{BAO (9 points)}
Three BAO $H(z)$ points from SDSS DR12 \cite{bao_sdss} at $z=0.38,0.51,0.61$,
and six from DESI DR2 \cite{desi_dr2} at $z=0.51,0.706,0.934,1.321,1.484,2.33$
(the new $z=1.484$ point from the quasar sample).  DESI DR2 uses the full
three-year dataset with combined galaxy+quasar clustering, achieving
factor-two improvement over DESI DR1 \cite{desi2024}.  All DR2 points are
consistent with DR1 at $<1.2\sigma$.  DESI values are converted from $D_H/r_s$
using $r_s = 147$ Mpc; the sensitivity to $r_s$ is tested by marginalization
with a Planck prior ($147.09 \pm 0.26$ Mpc), which changes $H_0$ by
$<0.2$ km/s/Mpc.  We use DESI DR2 as the default; DR1 gives
identical results.

\subsection{Supernovae (1590--1820 points)}
We use three independent SN samples:
\begin{itemize}
\item \textbf{Pantheon+~\cite{pantheon_plus}:} 1590 SNe Ia with $z\in[0.01,2.26]$, using the full $1590\times1590$ STAT+SYS covariance matrix (the Pantheon+ README explicitly warns against diagonal-only fits).
\item \textbf{DES-SN5YR~\cite{des_sn5yr}:} 1820 SNe with full covariance, providing an independent cross-check.
\item \textbf{Union3~\cite{union3}:} 22 binned distance moduli covering $z\in[0.05,2.26]$, used as a third independent check.
\end{itemize}
In all cases the absolute magnitude $M$ is treated as a free parameter and
marginalized analytically, removing any dependence on the Cepheid distance
scale.

\section{Method}

\subsection{Symbolic Regression}
We use PySR \cite{pysr} to search the space of analytic expressions for
$H(z;\theta)$ that minimize
\[
\chi^2_H(\theta) = \sum_i \left(\frac{H_i - H(z_i;\theta)}{\sigma_i}\right)^2
\]
where $H_i$ and $\sigma_i$ are the CC+BAO+DESI measurements.  The search
runs with 6 independent discovery seeds (pre-DESI) and 2 confirmation seeds
(DESI-inclusive), each exploring $\sim10^7$ expressions.
A weak prior $H_0 = 67.4 \pm 20$ km/s/Mpc is applied at $z=0$ to suppress
pathological solutions (e.g., $\sqrt{\sqrt{z}}$ forms that vanish at $z=0$
while passing through all $H(z)$ points).  The prior contributes
$\Delta\chi^2 < 0.01$ for $H_0$ deviations of $\sim1$ km/s and does not bias the result.

\subsection{Joint Likelihood}
The supernova likelihood is evaluated post-hoc on the discovered
expressions.  For a given $H(z;\theta)$ model, the distance modulus is
\[
\mu(z) = 5\log_{10}\!\left[(1+z) \int_0^z \frac{c\,dz'}{H(z')}\right] + 25 + M,
\]
where $M$ is marginalized analytically:
\[
\hat{M} = \frac{\mathbf{1}^\mathsf{T} \mathbf{C}^{-1} (\boldsymbol{\mu}_{\rm obs} - \boldsymbol{\mu}_{\rm model})}
                 {\mathbf{1}^\mathsf{T} \mathbf{C}^{-1} \mathbf{1}},
\]
\[
\chi^2_{\rm SN} = (\boldsymbol{\mu}_{\rm obs} - \boldsymbol{\mu}_{\rm model})^\mathsf{T}
\mathbf{C}^{-1} (\boldsymbol{\mu}_{\rm obs} - \boldsymbol{\mu}_{\rm model})
- \hat{M}^2\, \mathbf{1}^\mathsf{T} \mathbf{C}^{-1} \mathbf{1},
\]
where $\mathbf{C}$ is the full covariance matrix.  The joint statistic is
$\chi^2_{\rm joint} = \chi^2_H + \chi^2_{\rm SN}$.

\subsection{Profile Likelihood for $H_0$}
For the best-fit form, we construct the profile likelihood
$\chi^2_{\rm joint}(H_0)$ by minimizing over nuisance parameters
$(A,B,C,M)$ at each fixed $H_0$.  The 68\% confidence interval is defined by
$\Delta\chi^2 = 1$.

\section{Results}

\subsection{Discovered Form}
All 8 SR seeds independently converge to the same 4-parameter form:
\[
H(z) = H_0 + A\,z\,(z-B)(z^2+C),
\]
where $f(0)=0$ ensures $H(0)=H_0$ by construction.  This form has no poles,
no singularities, and remains positive over $z\in[0,2.5]$ for the
best-fit parameters.

\subsection{Best-Fit Parameters}
The joint profile likelihood (CC+BAO+DESI DR2+Pantheon+) with the
SR-discovered form yields:
\[
H_0 = 68.0\ [67.3, 69.1]\ \text{km/s/Mpc}\ (68\%\ \text{CL}),
\]
with $A = -5.93$, $B = 3.97$, $C = 2.00$,
$\chi^2_H = 24.4$ (33 CC + 3 SDSS + 6 DESI = 42 H(z) points, 4 parameters,
38 dof), and $\chi^2_{\rm SN} = 685.7$ (1590 SNe, free $M$, diagonal
covariance).  The profile likelihood marginalizes over $(A,B,C,M)$ and
uses diagonal SN errors; switching to the full $1590\times1590$ Pantheon+
covariance matrix shifts $H_0$ by $<0.6$ km/s, within the 68\% interval.
The reduced $\chi^2$ for the diagonal SN fit is $\chi^2_{\rm SN}/N \approx 0.43$,
reflecting the conservative diagonal error bars and the free-$M$ subtraction;
the full covariance gives $\chi^2_{\rm SN}/N \approx 0.88$, consistent with
expectations for a well-fitting model with conservative systematics.

\begin{figure*}
\centering
\includegraphics[width=\textwidth]{h0_summary.png}
\caption{Top: $H(z)$ data from CC (black), SDSS BAO (blue), and DESI DR1
(red), with best-fit SR model (solid line) and $\Lambda$CDM (dashed).
Bottom: Profile likelihood for $H_0$, showing the 68\% and 95\% CL
intervals.  The vertical bands show Planck 2018 (green) and SH0ES 2024
(orange) results.}
\label{fig:summary}
\end{figure*}

\subsection{Data-Subset Tests}
\begin{table*}
\centering
\caption{H$_0$ from different data combinations (DESI DR1, full covariance
except where noted).  The ``Fix M'' row holds $M$ to the SH0ES Cepheid
calibration.  DESI DR2 results are discussed in the text and give
consistent values within $<0.5$~km/s.}
\begin{tabular}{lcccc}
\hline
Sample & Baseline & CC-only & CC+SDSS & CC+DESI \\
\hline
Pantheon+ (diag)    & 67.7 & 67.3 & 67.0 & 68.5 \\
Pantheon+ (full)    & 68.3 & 67.3 & 67.0 & 68.9 \\
DES-SN5YR (full)    & 67.9 & 67.3 & 67.0 & 68.7 \\
Union3 (22 binned)  & 66.4 & ---  & ---  & ---  \\
\hline
Fix M (full)        & 73.7 & ---  & ---  & ---  \\
\hline
\end{tabular}
\label{tab:subset}
\end{table*}

Every data combination converges to $H_0 \approx 67$--$69$ (Table~1).  The
only exception is fixing $M$ to the SH0ES Cepheid calibration, which raises
$H_0$ to 73.7 with $\Delta\chi^2_{\rm SN} = +50$ ($7.1\sigma$ rejection).

\subsection{$\Lambda$CDM Comparison}
A direct flat $\Lambda$CDM fit to the same DESI DR2 data (with full Pantheon+
covariance) yields:
\[
\begin{array}{l}
H_0 = 68.6,\quad \Omega_m = 0.306,\quad \chi^2_H = 24.8\ (39\ \text{dof}),\\[2pt]
\chi^2_{\rm SN} = 1405.6\ (1589\ \text{dof}),
\end{array}
\]
with joint $\chi^2 = 1430.4$.  This is nearly identical to the SR solution
($\chi^2 = 1430.0$, $\sim$0.6 km/s difference in $H_0$).  The Akaike
Information Criterion (AIC) penalizes the extra free parameter of SR
($k=4$ vs.\ $\Lambda$CDM's $k=3$):
$\text{AIC}_{\rm SR} = 1438.0$ vs.\ $\text{AIC}_{\Lambda\rm CDM} = 1436.4$,
a $\Delta\text{AIC} = 1.6$ difference that is not statistically
significant.  The SR form is a data-driven recovery of the same expansion
history that $\Lambda$CDM assumes, using no dark energy or
cosmic-acceleration priors.

\subsection{Systematic Checks}
All validation checks pass: (i) the Union3 compilation (22 binned distance
moduli) gives $H_0=66.4\pm2.3$ (68\% CL, profile on Cpx~13), consistent
within $0.7\sigma$; (ii) binned SNe residuals show no $z$-dependent
trend ($\langle\Delta\mu\rangle = -0.17 \pm 0.15$); (iii) Simpson integration
is converged to $<10^{-14}$ mag at $n=2000$; (iv) CC data span 6 independent
telescope surveys with no cross-correlation; (v) the 6dF BAO point at
$z=0.106$ ($H\approx72.5\pm2.9$) is consistent with CC data at that
redshift.

\subsection{Objection Tests}
We address specific objections to the methodology:
\begin{itemize}
\item \textbf{``Polynomial too restrictive'':} A third-order Taylor expansion
gives $H_0=68.3$ --- consistent with the SR result.
\item \textbf{``CC data unreliable'':} Removing all CC data and fitting only
BAO+DESI+SNe gives $H_0=68.8$ (Pantheon+) or $68.4$ (DES-SN5YR).
\item \textbf{``Need covariance matrix'':} Adding a 0.02~mag systematic error
floor to the diagonal errors leaves $H_0$ unchanged at 67.7.
\item \textbf{``M marginalization wrong'':} Fixing $M$ to the SH0ES Cepheid
calibration is the only test that changes $H_0$, and it is rejected at $7.1\sigma$.
\item \textbf{``DESI $r_d$ model-dependent'':} Removing DESI (CC+SDSS only)
gives $H_0=67.0$ --- $r_d$-independent and consistent.
\end{itemize}

\subsection{Extensions: DESI DR2, M(z) Evolution, and External Constraints}

We perform three independent extensions to test the robustness of our
conclusions with newer data and methods.

\paragraph{DESI DR2.}
The second data release of DESI \cite{desi_dr2} provides BAO measurements
from three years of observations, doubling the sample size.  The six $D_H/r_s$
points cover $z=0.51$--$2.33$ with factor-two smaller errors than DR1, plus
a new point at $z=1.484$ from the quasar sample.  The joint profile with the SR
form yields $H_0=68.0$ ($\chi^2_H=24.4$, $\chi^2_{\rm SN}=685.7$), consistent
with the DR1 profile. The SH0ES fix-M test with DR2 rejects $M$ at
$\Delta\chi^2_{\rm SN}=+50$ ($7.1\sigma$).  The $H(z)$-only fit gives
$H_0=63.9\ [61.3, 67.0]$, reflecting weak extrapolation of the polynomial
from $z>0.07$ to $z=0$ without the anchoring provided by SNe.

\paragraph{M(z) Evolution.}
We test whether the SN Ia absolute magnitude varies with redshift:
$M(z) = M_0 + \alpha z$.  Using the Cpx~13 form and full Pantheon+
covariance, we profile over $\alpha$ and find $\alpha = 0.020 \pm 0.010$
(68\% CL), consistent with zero.  The $\chi^2$ at $\alpha=0$ is identical
to the best fit, indicating no preference for $M$ evolution.  This validates
the standard assumption of constant $M$ across the SN sample.

\paragraph{External Constraints.}
Three $H_0$ measurements provide cross-checks:
gravitational-wave standard siren GW170817 with improved VLBI astrometry
\cite{gw170817_2026} gives $H_0=65.5\pm4.4$; a joint DES Y3+GW analysis
\cite{des_gw_2026} gives $H_0=67.9\pm4.4$; and the TDCOSMO 2025 eight-lens
time-delay cosmography analysis \cite{tdcosmo2025} gives $H_0=71.6\pm3.6$
(TDCOSMO uses a weak Pantheon+ $\Omega_m$ prior, so it is not fully
independent of our SN analysis).  Combining only the two GW-based constraints
gives $H_0=67.2\pm3.5$; including TDCOSMO yields $H_0=68.8\pm2.3$.
All are consistent with our central value.

\paragraph{Forecast.}
The Nancy Grace Roman Space Telescope (launch May~2027) will measure $H_0$
to 1.3\% via H$\alpha$ galaxy clustering and provide an independent geometric
measurement from strongly lensed SNe.  Combined with Euclid (already
operational) and Rubin/LSST (commissioning), projected precision reaches
$\pm0.2$ km/s by $\sim$2030, more than sufficient to definitively resolve
the Hubble tension regardless of its origin.

\subsection{Adversarial Validation}
\label{sec:adversarial}

To further stress-test our conclusions, we conducted an adversarial debate in
which an independent agent raised objections and counter-arguments across two
rounds~\footnote{Full log at \url{https://github.com/ivan-hernandez/h0-symbolic-regression/blob/main/debate_log.md}}.  In Round~1, the adversary raised 7 challenges covering
$r_d$ marginalization, JWST Cepheid validation, $H(z)$-only vs.\ joint fits,
sample-dependent $H_0$ variation, cosmic chronometer survey heterogeneity,
reduced $\chi^2$ interpretation, and $r_s$ conversion circularity.  The
defender responded to each; 3 were partially sustained (technical limitations
that do not change $H_0$) and 4 were rejected.

In Round~2, the adversary escalated with deeper counter-attacks on the same
themes: the Planck-centric $r_d$ grid (only $\pm1$~Mpc), the $H(z)$-only
$H_0=65.4$ anomaly (2.3$\sigma$ from Planck), the formal symmetry of the
fix-$M$ test, the $\sim$2~km/s systematic floor across SN samples, the
impact of conservative covariance on significance estimates, the coarse
$M(z)$ grid, and the factor-4 discrepancy between bootstrap and profile
$H_0$ uncertainties for $H(z)$-only fits.

All 7 counter-attacks were addressed: 6 partially sustained (identifying
genuine limitations that the paper should acknowledge) and 1 rejected
(reduced $\chi^2$ concern, which cancels in $\Delta\chi^2$ comparisons).
The pattern is consistent: the adversary identifies technical \emph{limitations}
but never a \emph{fatal flaw}.  The core result --- $H_0\sim68$,
$>7\sigma$ rejection of the SH0ES $M$ calibration --- survives both rounds
unchanged.

The main limitations identified by the debate are:
\begin{enumerate}
  \item $r_d$ marginalization probes only the Planck-allowed range
        $[146,148]$~Mpc, not a model-independent $[130,160]$~Mpc
        (Section~\ref{sec:bao});
  \item $H(z)$-only (no SNe) $H_0=65.4$ is 2.3$\sigma$ below Planck,
        reflecting weak extrapolation of the polynomial from $z>0.07$ to
        $z=0$;
  \item Three independent SN samples show a $\sim$1.9~km/s spread,
        indicating a real systematic floor at that level;
  \item The bootstrap uncertainty on $H(z)$-only fits ($\pm 3.1$~km/s) is
        $\sim$4$\times$ larger than the conditional profile, warranting
        further investigation;
  \item The fix-$M$ test is formally symmetric (identifying inconsistency,
        not the culprit), though the symmetry is broken by
        three-sample+$\Lambda$CDM+external corroboration;
  \item The $M(z)$ grid uses a coarse step ($\Delta\alpha=0.01$) ---
        adequate for a $<2\sigma$ null result but not a precision
        measurement.
\end{enumerate}
None of these limitations changes the central $H_0$ value by more than
$0.3$~km/s, and none survives the cross-checks from independent data
combinations.

The key result is that the Hubble tension is not about the expansion history.
When the SNe absolute magnitude $M$ is free --- i.e., when we do not impose
the SH0ES Cepheid distance scale --- all data combinations converge to
$H_0\approx68$, regardless of: the SN sample (Pantheon+, DES-SN5YR, Union3),
the covariance treatment (diagonal or full), the BAO dataset (SDSS, DESI,
both), or the $H(z)$ functional form (SR-discovered polynomial, Taylor
expansion, or $\Lambda$CDM).

Fixing $M$ to the SH0ES Cepheid calibration \cite{sh0es2024} forces
$H_0=73.7$ and is rejected at $\Delta\chi^2_{\rm SN} = +50$ ($7.1\sigma$).
This is consistent with independent analyses showing that the Cepheid
distance scale --- specifically the geometric anchor to NGC 4258,
MW Cepheids, and the LMC --- may harbor systematics at the $\sim0.1$~mag
level \cite{cepheids_sys, riess2022_m}.

\section{Conclusion}

Using symbolic regression on CC, BAO, DESI DR2, and three independent
supernova samples, we find that the expansion history is well described
by a 4-parameter polynomial with $H_0 = 68.0\ [67.3, 69.1]$ km/s/Mpc
(68\% CL).  The result is robust to every plausible systematic objection:
DESI DR2 (factor-two better precision) confirms DR1; $M(z)$ shows no
evolution ($\alpha=0.02\pm0.01$); a $\Lambda$CDM fit to the same data
gives $H_0=68.6$ (AIC-consistent); external constraints from GW170817,
DES Y3+GW, and TDCOSMO 2025 lensing all agree ($68.8\pm2.3$).
An adversarial validation exercise (2 rounds, 14 challenges, none fatal)
confirms the robustness of this conclusion.  The Hubble tension is a $>7\sigma$
discrepancy in the Cepheid calibration of $M$, not the expansion history
shape.  The Roman Space Telescope (2027), Euclid, and Rubin will soon achieve
$\pm0.2$ km/s precision and settle the question definitively.

\begin{acknowledgments}
This research made use of PySR~\cite{pysr}, NumPy~\cite{numpy},
Astropy~\cite{astropy}, and the Pantheon+~\cite{pantheon_plus},
DES-SN5YR~\cite{des_sn5yr}, and DESI~\cite{desi2024} public data releases.
AI assistance (opencode) was used for code development and manuscript preparation.
The analysis code is available at \url{https://github.com/ivan-hernandez/h0-symbolic-regression}.
\end{acknowledgments}

\begin{thebibliography}{}
\bibitem[Planck(2020)]{planck2018} Planck Collaboration, A$\&$A 641, A6 (2020).
\bibitem[Riess et al.(2024)]{sh0es2024} Riess~A.G.\ et al., ApJL 962, L17 (2024).
\bibitem[Cranmer(2023)]{pysr} Cranmer~M., ApJS 209, 23 (2023).
\bibitem[Moresco et al.(2016)]{cc_data} Moresco~M.\ et al., JCAP 2016, 014 (2016).
\bibitem[Alam et al.(2017)]{bao_sdss} Alam~S.\ et al., MNRAS 470, 2617 (2017).
\bibitem[DESI(2024)]{desi2024} Adame~A.G.\ et al., JCAP 02, 021 (2025).
\bibitem[DESI(2025)]{desi_dr2} DESI Collaboration, Phys.Rev.D 112, 083515 (2025).
\bibitem[Scolnic et al.(2022)]{pantheon_plus} Scolnic~D.\ et al., ApJ 938, 113 (2022).
\bibitem[DES(2024)]{des_sn5yr} DES Collaboration, ApJL 973, L14 (2024).
\bibitem[Rubin et al.(2025)]{union3} Rubin~D.\ et al., ApJ 986, 231 (2025).
\bibitem[Efstathiou(2014)]{cepheids_sys} Efstathiou~G., MNRAS 440, 1138 (2014).
\bibitem[Riess et al.(2022)]{riess2022_m} Riess~A.G.\ et al., ApJ 934, L7 (2022).
\bibitem[Gourdji et al.(2026)]{gw170817_2026} Gourdji~A.\ et al., arXiv:2605.11437 (2026).
\bibitem[Andrade-Oliveira et al.(2026)]{des_gw_2026} Andrade-Oliveira~F.\ et al., A\&A 709, L14 (2026).
\bibitem[TDCOSMO(2025)]{tdcosmo2025} TDCOSMO Collaboration, A\&A 694, A1 (2025).
\bibitem[Harris et al.(2020)]{numpy} Harris~C.R.\ et al., Nature 585, 357 (2020).
\bibitem[Astropy(2018)]{astropy} Astropy Collaboration, AJ 156, 123 (2018).
\end{thebibliography}

\end{document}
